\chapter{章节标题}

\section{本章要解决的问题}

从一个具体、可运行的 C/C++ 问题开始。不要先列定义。

\begin{keyidea}
用一段话写本章最重要的思想。
\end{keyidea}

\section{收束目标}

本章收束时，应能：

\begin{itemize}
  \item 解释某个机制为什么存在。
  \item 画出底层数据流或控制流。
  \item 用命令观察这个机制。
  \item 写代码验证一个现象。
  \item 用性能或正确性证据支持结论。
\end{itemize}

\section{第一层：源码层抽象}

用 C/C++ 例子讲清表层模型。

\section{第二层：编译器或操作系统看到的模型}

解释该抽象如何被降低、检查、优化或管理。

\section{第三层：硬件或运行时实现模型}

讲底层如何大致实现，不要求芯片级细节，但要足够解释现象。

\section{贯穿例子：从源码到证据}

\subsection{源码}

\begin{lstlisting}[language=C++]
// minimal example
\end{lstlisting}

\subsection{命令}

\begin{lstlisting}[language=bash]
# commands
\end{lstlisting}

\subsection{汇编、IR 或工具输出}

\begin{lstlisting}
# selected output
\end{lstlisting}

\subsection{解释}

逐行解释输出。说明每个观察点支持什么结论。

\section{边界和反例}

至少写两个容易误解的反例。

\section{实验}

\begin{labbox}
写实验目的、代码、命令、观察点、预期结果和异常结果解释。
\end{labbox}

\section{习题}

\begin{exercisebox}
至少包含解释题、代码题、工具观察题和证据记录题。
\end{exercisebox}

\section{常见误区}

\begin{warningbox}
\begin{itemize}
  \item 误区 1。
  \item 误区 2。
  \item 误区 3。
\end{itemize}
\end{warningbox}

\section{验收标准}

\begin{itemize}
  \item 能画图。
  \item 能解释。
  \item 能运行命令。
  \item 能写实验报告。
\end{itemize}
